%; whizzy paragraph -pdf xpdf -latex ./whizzypdfptex.sh
%; whizzy-paragraph "^\\\\begin{frame}\\|\\\\emtext"
% latex beamer presentation.
% platex, latex-beamer でコンパイルすることを想定。 

%     Tokyo Debian Meeting resources
%     Copyright (C) 2012 Junichi Uekawa
%     Copyright (C) 2016 Nobuhiro Iwamatsu

%     This program is free software; you can redistribute it and/or modify
%     it under the terms of the GNU General Public License as published by
%     the Free Software Foundation; either version 2 of the License, or
%     (at your option) any later version.

%     This program is distributed in the hope that it will be useful,
%     but WITHOUT ANY WARRANTY; without even the implied warreanty of
%     MERCHANTABILITY or FITNESS FOR A PARTICULAR PURPOSE.  See the
%     GNU General Public License for more details.

%     You should have received a copy of the GNU General Public License
%     along with this program; if not, write to the Free Software
%     Foundation, Inc., 51 Franklin St, Fifth Floor, Boston, MA  02110-1301 USA

\documentclass[cjk,dvipdfmx,12pt]{beamer}
\usetheme{Tokyo}
\usepackage{monthlypresentation}

%  preview (shell-command (concat "evince " (replace-regexp-in-string "tex$" "pdf"(buffer-file-name)) "&")) 
%  presentation (shell-command (concat "xpdf -fullscreen " (replace-regexp-in-string "tex$" "pdf"(buffer-file-name)) "&"))
%  presentation (shell-command (concat "evince " (replace-regexp-in-string "tex$" "pdf"(buffer-file-name)) "&"))

%http://www.naney.org/diki/dk/hyperref.html
%日本語EUC系環境の時
\AtBeginDvi{\special{pdf:tounicode EUC-UCS2}}
%シフトJIS系環境の時
%\AtBeginDvi{\special{pdf:tounicode 90ms-RKSJ-UCS2}}

\newenvironment{commandlinesmall}%
{\VerbatimEnvironment
  \begin{Sbox}\begin{minipage}{1.0\hsize}\begin{fontsize}{8}{8} \begin{BVerbatim}}%
{\end{BVerbatim}\end{fontsize}\end{minipage}\end{Sbox}
  \setlength{\fboxsep}{8pt}
% start on a new paragraph

\vspace{6pt}% skip before
\fcolorbox{dancerdarkblue}{dancerlightblue}{\TheSbox}

\vspace{6pt}% skip after
}
%end of commandlinesmall

\title{Mini Debian Conference Japan 2016}
\subtitle{The Japan Technical Jamboree 59}
\author{Debian Project / Debian JP Project \\岩松 信洋}
\date{2016年12月2日}
\logo{\includegraphics[width=8cm]{image200607/openlogo-light.eps}}

\begin{document}

\begin{frame}
\titlepage{}
\end{frame}

\begin{frame}
\begin{center}
\Huge{Mini Debian Conference Japan 2016 開催のお知らせ}
\end{center}
\end{frame}

\begin{frame}
\begin{center}
\Huge{\url{miniconf.debian.or.jp}}
\end{center}
\end{frame}

\begin{frame}{開催概要}
  \begin{itemize}
   \item 日時: 12/10 10:30 - 
   \item 場所: サイボウズ株式会社 / 東京都中央区日本橋2-7-1 東京日本橋タワー
   \item URL: \url{http://miniconf.debian.or.jp}
   \item 懇親会、GPGキーサインパーティあり。
  \end{itemize}
\end{frame}

\begin{frame}{本カンファレンスの目的}
  \begin{itemize}
   \item 日本でのDebian / FLOSS 開発活性化
   \item Debian 開発者とDebian/Ubuntuユーザとの情報交換の場
   \begin{itemize}
   	\item 毎月関東と関西でDebianの勉強会も行っています。
   \end{itemize}
   \item Debian Conference を日本で開催するための周知および調査活動
  \end{itemize}
\end{frame}

\begin{frame}{発表（一部）}
  \begin{itemize}
   \item Raspbian と OpenCV の話
   \item Go / emacs ソフトウェアパッケージング
   \item Debianソースコードを用いた組込み向けLinux環境作成 （小林さん）
   \item OSS ボーカロイド
   \item Firefox カスタマイズ （クリアコード/林さん）
   \item Debian Conference in Taiwan について
   \item etc.
  \end{itemize}
\end{frame}

\begin{frame}{開催概要}
  \begin{itemize}
   \item 日時: 12/10 10:30 - 
   \item 場所: サイボウズ株式会社 / 東京都中央区日本橋2-7-1 東京日本橋タワー
   \item URL: \url{http://miniconf.debian.or.jp}
   \item 懇親会、GPGキーサインパーティあり。
  \end{itemize}
\end{frame}

\end{document}

;;; Local Variables: ***
;;; outline-regexp: "\\([ 	]*\\\\\\(documentstyle\\|documentclass\\|emtext\\|section\\|begin{frame}\\)\\*?[ 	]*[[{]\\|[]+\\)" ***
;;; End: ***
